\partsection{B-01}{API tests on the old and new rule}
\label{ev:redgreen}
On 24 November, the team captured RED on \texttt{214d1d5} with the revised tests from \texttt{405c906} (service) and \texttt{6cabc3e} (HTTP). GREEN ran those same tests on \texttt{b51c058}. Exact test start times were unavailable.

\textbf{Setup and assertions.}
In \texttt{apps/api/internal/claim/act\_test.go}, F04-TEST-06 creates a PostgreSQL-backed submitted claim per case. Its assigned reviewer calls \texttt{Act} with the current version. A failed decision must give exactly one \texttt{comment}/BR-04 error and leave the claim submitted at its original version, with no decision time and one history entry. Success stores the trimmed reason.

In \texttt{apps/api/claims\_review\_test.go}, F04-TEST-07 logs in as the assigned reviewer on a newly submitted claim. Each 422 must return the expected message and leave the database state equal to the initial snapshot. Each accepted length uses a fresh claim. Selected test body from the 422 loop:
\begin{record}
recorder := decide(t, router, token, id, tc.action, "\"0\"", tc.body)
decodeProblem(t, recorder, http.StatusUnprocessableEntity, "VALIDATION_FAILED")

want := []fieldError{{Field: "comment", Rule: "BR-04", Message: tc.message}}
if got := problemErrors(t, recorder); !slices.Equal(got, want) {
    t.Fatalf("errors = %+v, want %+v", got, want)
}

if after := readState(t, pool, id); after != before {
    t.Fatalf("stored = %+v after the refused decision, want %+v", after, before)
}
\end{record}

\textbf{Commands and captured output.}
Both commands ran from \texttt{apps/api} on RED and GREEN.
\begin{record}
go test -v -count=1 ./internal/claim -run '^TestF04TEST06CommentRuleBR04$'
go test -v -count=1 . -run '^TestF04TEST07CommentRuleIs422WithTheFieldError$'
\end{record}
\begin{record}
RED (exit 1 for both commands)
=== RUN   TestF04TEST06CommentRuleBR04/reject_with_9_characters
    act_test.go:87: Act error = <nil>, want a ValidationError with one field
--- FAIL: TestF04TEST06CommentRuleBR04
=== RUN   TestF04TEST07CommentRuleIs422WithTheFieldError/reject_with_9_characters
    claims_review_test.go:483: status = 200, want 422
--- FAIL: TestF04TEST07CommentRuleIs422WithTheFieldError

GREEN (exit 0 for both commands)
--- PASS: TestF04TEST06CommentRuleBR04
    --- PASS: TestF04TEST06CommentRuleBR04/reject_with_9_characters
    --- PASS: TestF04TEST06CommentRuleBR04/reject_with_10_characters
--- PASS: TestF04TEST07CommentRuleIs422WithTheFieldError
    --- PASS: TestF04TEST07CommentRuleIs422WithTheFieldError/reject_with_9_characters
    --- PASS: TestF04TEST07CommentRuleIs422WithTheFieldError/reject_with_10_characters
\end{record}
RED had 11 pass / 5 fail in the module and 4 pass / 5 fail at HTTP. Old code accepted the 9-character input and used the old blank-reason message. GREEN had 16 module and 9 HTTP passes, 0 failures and 0 skips. GREEN also accepted 500, rejected 501, and preserved state on each 422.

\clearpage
\partsection{B-02}{Dialog and seeded test evidence}
\label{ev:web}
\textbf{F04-TEST-21 (component).}
With a submitted claim and a mocked decision endpoint, the test types a padded 9-character reason and checks the message, disabled state and absence of a POST. It checks Unicode 9/10 cases with edge whitespace, and an empty field. A separate mock returns 422 for “No receipt” (10) to verify API field errors and input retention.
\begin{record}
bun run test run src/features/reviews/__test__/decision-panel.test.tsx
\end{record}
Ran in \texttt{apps/web} on 24 November. At \texttt{6cabc3e}, the short-reason test expected \texttt{ariaDisabled} to be \texttt{"true"} and received \texttt{null}. Counts were 13 passed, 1 failed, 0 skipped. At \texttt{b51c058}, all 14 passed with 0 failed and 0 skipped. The teammate's rerun matched these counts.

\textbf{F04-TEST-24 (Chromium E2E).}
The requester submits a new claim and the assigned reviewer opens it from the queue. The test checks empty and 9-character disabling, submits a valid reason, and opens the claim again as the requester to confirm the transition and stored reason. The paired F04-TEST-23 checks approval and requester history.
\begin{record}
bun run e2e e2e/f04-review-decision.spec.ts
\end{record}
Ran from \texttt{apps/web} on 24 November against local PostgreSQL 17 with seven seeded users and six sample claims. On the RED dialog and spec from \texttt{6cabc3e} with the updated API, the empty-reason check showed \texttt{Expected: disabled} and \texttt{Received: enabled}. This isolates the old dialog. The record gives one failing case. The full RED run totals are unavailable. On \texttt{0226536} with API \texttt{b51c058}, both journeys passed (2 passed, 0 failed, 0 skipped). The follow-up narrowed the Playwright locator to the reason field's accessible description.

\partsection{B-03}{OpenAPI contract and implementation diff}
\label{ev:diff}
The schema edit in \texttt{ea1537b} changed the reject endpoint description and shared \texttt{ClaimDecision.comment} description to 10--500 after trimming. The service returns a 422 with the field error. The shared schema gained metadata while approve kept its optional comment:
\begin{record}
+          x-min-length: 10
           x-max-length: 500
\end{record}
Selected actual lines from the \texttt{b51c058} diff in \texttt{apps/api/internal/claim/rules.go}. The existing trimmed rune count applies before these branches.
\begin{record}
-   MinReasonLength  = 5
+   MinReasonLength  = 10
    MaxCommentLength = 500

    case action == ActionReject && length == 0:
-       message = "Enter a reason (5 to 500 characters)."
+       message = fmt.Sprintf("Enter a reason (%d to %d characters).", MinReasonLength, MaxCommentLength)
    case action == ActionReject && length < MinReasonLength:
-       message = "Reason must be at least 5 characters."
+       message = fmt.Sprintf("Reason must be at least %d characters.", MinReasonLength)
\end{record}

\clearpage
\textbf{B-03 continued.}
The web generator propagates \texttt{x-min-length} into Zod metadata, and \texttt{contract.ts} reads it. Selected actual lines from \texttt{decision-panel.tsx} at \texttt{b51c058} show the reject-only gate and field error:
\begin{record}
+const commentSchema = zClaimDecision.shape.comment.unwrap();
+const reasonMinimum = contractMinLength(commentSchema);
+const commentLimit = contractMaxLength(commentSchema);

+  const reasonTooShort =
+    decision === "reject" &&
+    (comment
+      .replaceAll(
+        /^\p{White_Space}+|(?<!\p{White_Space})\p{White_Space}+$/gu,
+        ""
+      )
+      .match(/./gsu)?.length ?? 0) < reasonMinimum;
+
+  const minimumError =
+    reasonTooShort && comment !== ""
+      ? `Reason must be at least ${reasonMinimum} characters.`
+      : undefined;

+  const commentError = minimumError ?? serverCommentError;
-                isDisabled={send.isPending}
+                isDisabled={send.isPending || reasonTooShort}
\end{record}
Generated outputs changed in \texttt{api.gen.go}, \texttt{types.gen.ts}, \texttt{sdk.gen.ts}, \texttt{react-query.gen.ts} and \texttt{zod.gen.ts}. Generation left a clean diff. The implementation uses the existing write path, permissions, \texttt{If-Match} checks and transaction. No migration or history rewrite was added.

\partsection{B-04}{Regression selection and results}
\label{ev:regression}
The 24 November regression set exercises F-04 dependents. F04-TEST-05 checks the reason in requester history. F04-TEST-08 to 13 cover owner and reviewer access. F04-TEST-14 to 18 cover status, stale and concurrent decisions, rollback and append-only history. F04-TEST-23 covers approval. The web suite also covers the review queue, session and decision errors. Revised expectations are listed in A-03 on p.~\pageref{ev:contract-review}.
\begin{longtable}{@{}R{10.4cm}R{5.7cm}@{}}
\toprule
\textbf{Command and version} & \textbf{Actual result} \\ \midrule \endhead
\texttt{go test -v -count=1 . ./internal/claim -run '\string^TestF04TEST'} at \texttt{b51c058} and \texttt{go test -v -count=1 ./migrations -run '\string^TestF04TEST'} at \texttt{0226536}, both in \texttt{apps/api}. &
19 named F-04 tests passed. No failures or skips reported for this selection. \\
\texttt{bun run test:coverage} in \texttt{apps/api} at \texttt{b51c058}. &
All thresholds passed, 100\% (785/785). The local record gives coverage only. See B-05 for release counts. \\
\texttt{bun run test:coverage --maxWorkers=2 --retry=2} in \texttt{apps/web} at \texttt{0226536}. &
139 passed in 17 files. Statement, branch, function and line coverage were 100\%. Earlier concurrent F01-TEST-20 runs hit login timeouts. \\
Local seeded F-04 E2E at \texttt{0226536} (command and setup in B-02). &
2 passed, 0 failed, 0 skipped. \\
\texttt{bun run generate} and \texttt{bun run vuln} at \texttt{b51c058}. &
Generation left no diff. API scan found no vulnerabilities. \\
\bottomrule
\end{longtable}

\clearpage
\partsection{B-05}{Review, staging and released main}
\label{ev:release}
\begin{longtable}{@{}R{3.8cm}R{12.3cm}@{}}
\toprule
\textbf{24 Nov 2026 (UTC)} & \textbf{Recorded action and review evidence} \\ \midrule \endhead
16:16--16:18 & Contract \prlink{101} approved, then squashed to \texttt{lab-7} as \texttt{124cf9a}. See A-03. \\
16:20--16:28 & OpenAPI contract commit \texttt{ea1537b} came first. Test commits \texttt{405c906} and \texttt{6cabc3e} preceded code \texttt{b51c058}. Implementation proceeded at the maintainer's direction while OpenAPI teammate review was pending. \\
17:16:23 & On \texttt{3e4b572}, \texttt{Kiatisakk} \href{\repo/pull/105\#pullrequestreview-6043999986}{confirmed the OpenAPI diff matched F-04}. The reviewer checked metadata generation, API RED/GREEN, and 50 focused web tests. Requested changes were to record the late review and list RED/GREEN versions per test. The reviewer relied on CI and the PR for full gates and E2E. \\
17:20--17:31 & \texttt{4993321} documented the review. \texttt{9b1fc9c} listed each test's versions. \texttt{Kiatisakk} approved at 17:28:21 after a docs-only follow-up and confirmed the component counts. \prlink{105} squashed to \texttt{lab-7} as \texttt{8b660a1} at 17:31:06. \\
17:47--17:48 & \texttt{Kiatisakk} approved \prlink{106} at 17:47:07 after checking the branch contents, generated rule, and required CI. The reviewer left the four authenticated checks for after deploy. \prlink{106} merged to \texttt{main} at 17:48:05. \issuelink{99} closed as completed at 17:48:07. \\
17:53:02 & Semantic-release published \texttt{v0.1.0.0}. Remote \texttt{main} and that tag both resolved to \releaseSHA{}. \\
\bottomrule
\end{longtable}
The release review calls \texttt{871d373} an approved head. Git identifies it as the tree ID shared by the actual approved \prlink{105} head \texttt{9b1fc9c}, staging squash \texttt{8b660a1} and release merge \texttt{c32e645}. The release's first parent is \texttt{214d1d5} and its second parent is \texttt{8b660a1}. This preserves the \texttt{feat!:} release signal. The 19 changed files cover L7/01 only. The separate Sprint 2 contract PR \#100 and release path \#102 are outside this change.

\textbf{CI.}
The five required jobs (\texttt{api}, \texttt{test}, \texttt{code\_quality}, \texttt{image}, \texttt{build}) passed on \prlink{101} run \runlink{45434761904}, \prlink{105} run \runlink{45447692201}, and \prlink{106} run \runlink{45449479661}. The release job was skipped on PRs. On main run \runlink{45452555188}, all five jobs and the release job succeeded. The release CI ran at \texttt{c32e645}, from 17:48:08 to 17:54:08.

Selected main CI output, with ANSI formatting and runner prefixes removed:
\begin{record}
api   17:52:19  DONE 300 tests in 57.660s
api   17:52:28  Total test coverage: 100% (785/785)
api   17:52:35  No vulnerabilities found.
test  17:48:50  Test Files  17 passed (17)
test  17:48:50  Tests       139 passed (139)
\end{record}
The Go job ran \texttt{gotestsum} with \texttt{-race -coverpkg=./... -coverprofile=coverage.out ./...} in \texttt{apps/api}. Its JUnit artifact (14949441476) confirms 300 passed, 0 failed, 0 errors and 0 skips. The web job ran \texttt{bun run test:coverage} in \texttt{apps/web}, with text-summary, Cobertura, and JUnit outputs. It recorded 0 failures and 0 skips. Local Playwright was run separately, per \prlink{96}.

\clearpage
\partsection{B-06}{Cloud Run deployment and public smoke}
\label{ev:deploy}
Deploy demo \runlink{45453647851} succeeded on 24 November from 17:54:12 to 17:56:06 after the main CI pass. It built from \releaseSHA{}, ran the existing migrate-and-seed job, and deployed by digest to \revision{}. The deployment step reported 100\% traffic to that revision. The digest was:
\begin{record}
sha256:4648f547d5250c04c9240a2b8b4ffe18a3fb4c0cf90501f949d13263f55dcb7e
\end{record}

The workflow ran \texttt{bash deploy/smoke.sh} with the public service URL and expected full SHA. Selected output from 17:55:49--17:55:59:
\begin{record}
Service [claimflow] revision [claimflow-00016-zck] has been deployed and is serving 100 percent of traffic.
PASS /api/v1/health reports c32e6455724a0b7be3a5ae8a50889ea50fd4ddbc
PASS /api/v1/ready 200
PASS / serves the web build
PASS unknown /api route is a JSON 404
\end{record}
The report session reran this read-only smoke from the release checkout against the public URL at about 18:23 UTC. All four checks passed again and the script exited 0.

\partsection{B-07}{Authenticated post-deployment checks}
\label{ev:live}
Automated HTTP probe run on demand against the demo on 24 November, 18:25:50--18:26:14 UTC (25 November, 01:25--01:26 ICT). The input was one fresh synthetic claim, \texttt{CLM-2026-000023}, created with a 250.00 THB travel item. Three role-specific logins returned 200. Creation returned 201. The assigned reviewer read \texttt{SUBMITTED}, ETag \texttt{"0"}, and history \texttt{[SUBMIT]}.

Each attempt called \texttt{POST /api/v1/claims/\{id\}/reject} with \texttt{If-Match: "0"}. After each refused request, an owner read compared status, version, decision time and history with the initial record. Public health reported \texttt{c32e645} before and after the verified run on \revision{}.

\begin{longtable}{@{}R{4.2cm}R{4.8cm}R{7.1cm}@{}}
\toprule
\textbf{Actor and input} & \textbf{Expected result} & \textbf{Actual result at 18:26:13 UTC} \\ \midrule \endhead
Owner \texttt{req.anan}, \texttt{abcdefghij} (10) & 403, no claim data or state change. & \textbf{PASS}. 403 \texttt{FORBIDDEN}, problem fields only, unchanged. \\
Unassigned reviewer \texttt{rev.somchai}, \texttt{abcdefghij} (10) & 403, no claim data or state change. & \textbf{PASS}. 403 \texttt{FORBIDDEN}, problem fields only, unchanged. \\
Assigned reviewer \texttt{rev.malee}, \texttt{abcdefghi} (9) & 422 on \texttt{comment}, BR-04. No state change. & \textbf{PASS}. 422 \texttt{VALIDATION\_FAILED}, field \texttt{comment}, rule BR-04, “Reason must be at least 10 characters.” ETag stayed \texttt{"0"}. \\
Assigned reviewer \texttt{rev.malee}, \texttt{abcdefghij} (10) & 200, \texttt{REJECTED}, version +1, one new history entry. & \textbf{PASS}. 200. Owner reread 200, ETag \texttt{"1"}, \texttt{REJECTED}, history \texttt{[SUBMIT, REJECT]} with reason \texttt{abcdefghij} and assigned reviewer as actor. \\
\bottomrule
\end{longtable}
The first attempt stopped after the owner 403 because the probe treated the standard problem \texttt{instance} request path as returned claim data. The corrected assertion checked the \texttt{instance} echo separately and confirmed the remaining fields contain no claim data. The rerun reused the same claim, finished with 4 passes, 0 failures and 0 skips, and ended all three sessions with 204. The synthetic rejection remains in append-only demo history.
